\begin{lemma}{Absch�tzung eines Gammafunktionsquotienten}
             {AbschaetzungEinesGammafunktionsquotienten}
Seien $\Gamma: ]0, \infty[ \rightarrow \R, 
             t \mapsto \int_0^\infty x^{t-1}\exp(-x)dx$ und
$a : \N \rightarrow \R, 
  n \mapsto \sqrt 2 \frac{\Gamma\left(\frac{n+1}2 \right)}
{\Gamma\left(\frac n 2 \right)}$
(nach \refclaim{ElementareEigenschaftenDerGammafunktion}
{ElementareEigenschaftenDerGammafunktion1} ist $a$ wohldefiniert).\\
Dann gilt:
\begin{enumerate}[(1)]
\item\label{AbschaetzungEinesGammafunktionsquotienten1} 
 $\forall \, n \in \N: a_n\cdot a_{n+1} = n$
\item\label{AbschaetzungEinesGammafunktionsquotienten2}
 $\forall \, n \in \N: a_n \le \sqrt n$
\item\label{AbschaetzungEinesGammafunktionsquotienten3}
 $\forall \, n \in \N: a_n \ge \frac{n}{\sqrt{n+1}}$
 \item\label{AbschaetzungEinesGammafunktionsquotienten4}
 $\forall\, n \in \N: \sqrt{\frac 2 \pi} \sqrt n \le a_n$
\end{enumerate}
\end{lemma}
\begin{proof}
Nach \refclaim{ElementareEigenschaftenDerGammafunktion}
{ElementareEigenschaftenDerGammafunktion1} ist f�r alle $n \in \N$: 
$\Gamma\left(\frac n 2 \right) > 0$ und $\Gamma(\frac{n+1}2)> 0$.\\
Daher ist f�r alle $n \in \N$: $a_n > 0$.\\
(1) Sei also $n \in \N$. Dann gilt:
\begin{align*}
a_n a_{n+1} &= 
\sqrt 2 \frac{\Gamma\left(\frac{n+1}2 \right)} {\Gamma\left(\frac n 2 \right)}
\cdot
\sqrt 2\frac{\Gamma\left(\frac{n+2}2 \right)}{\Gamma\left(\frac{n+1}2\right)}
= 2 \frac{\Gamma\left(\frac{n+2}2 \right)}{\Gamma\left(\frac n 2 \right)}
= 2 \frac{\Gamma\left(\frac n 2 + 1\right)}{\Gamma\left(\frac n 2 \right)}
\overset{\refclaim{ElementareEigenschaftenDerGammafunktion}
{ElementareEigenschaftenDerGammafunktion2}}= 
2\frac n2 \frac{\Gamma\left(\frac n2\right)}{\Gamma\left(\frac n2\right)}\\
&=n\text.
\end{align*}
(2)/(4) Sei $n \in \N$.\\
Dann gibt es nach  \refclaim{ExistenzFVerteilterZufallsgr�ssen}
{ExistenzFVerteilterZufallsgr�ssen2}  einen Wahrscheinlichkeitsraum 
$(\Omega', \mfa', \Prim')$ und ein
$\gamma: \Omega' \rightarrow \R^n$ unabh�ngig und $N(0, 1)$--verteilt.\\
Nach \ref{ErwartungswertDer2NormEinesStandardnormalverteiltenProzesses} gilt
dann $\E\left(\norm\cdot_{2, n} \circ \gamma \right) = a(n)$.\\
Da $\gamma$ unabh�ngig und $N(0,1)$--verteilt ist, ist es auch
unkorreliert, zentriert und $L_2-$normiert, sodass mit
\refclaim{ElementareEigenschaftenQuadratintegrierbarerAbbildungen}
{ElementareEigenschaftenQuadratintegrierbarerAbbildungen4} auch
$a(n) = \E\left(\norm\cdot_{2, n} \circ \gamma \right) \le \sqrt n$ gilt.\\
Au�erdem gilt nach \refclaim{AequivalenzkonstantenFuerPNormen}
{AequivalenzkonstantenFuerPNormen1}: $\frac 1 {\sqrt n} \norm\cdot_{1,n}
\le \norm\cdot_{2, n}$.\hfill$(\ast)$\\
Also gilt $0\le\frac 1{\sqrt n}\norm\cdot_{1,n}\circ\gamma\le\norm\cdot_{2,n}
\circ \gamma$ und nach
\ref{ErwartungswertDer2NormEinesStandardnormalverteiltenProzesses} ist
$\norm\cdot_{2, n} \circ \gamma$
integrierbar, also nach \refclaim{Integrierbarkeitskriterien}
{Integrierbarkeitskriterien1} auch $\norm\cdot_{1,n} \circ \gamma$.\\
Ferner ist $\E(\abs{\gamma_1}) = \sqrt{\nicefrac 2 \pi}$ und da die
Komponentenfunktionen von $\gamma$ jeweils
$N(0,1)$--verteilt sind, gilt: $\E(\norm\cdot_{1,n} \circ \gamma) 
= \sqrt{\nicefrac 2 \pi}\cdot n$.\\
Somit gilt nach \refclaim{Integrierbarkeitskriterien}
{Integrierbarkeitskriterien1}
\begin{align*}\E\left(\norm\cdot_{2,n} \circ \gamma \right)
 \overset{(\ast)}\ge \E \left(\frac 1 {\sqrt n} \norm\cdot_{1,n}
\circ \gamma \right) = \frac 1 {\sqrt n} \E \left(\norm\cdot_{1,n}\circ \gamma
\right)
= \frac 1 {\sqrt n} \sqrt{\frac 2 \pi} n = \sqrt{\frac 2 \pi} \sqrt n\text.
\end{align*}
(3) Sei $n \in \N$. 
Dann gilt: 
$a_n \overset{\eqref{AbschaetzungEinesGammafunktionsquotienten1}}= 
\frac n {a_{n+1}}
\overset{\eqref{AbschaetzungEinesGammafunktionsquotienten2}}\ge 
\frac n {\sqrt{n+1}}$.\\
\end{proof}